\chapter{FlashAttention — same math, much better hardware story}
\label{ch:flashattention}

\newthought{The first thing to try} when faced with a costly algorithm is
to see if you have been computing it stupidly. \citet{dao2022flashattention}
asked exactly this question of dense attention and discovered that, on
modern GPUs, the standard implementation was leaving an enormous amount of
performance on the table — not because the algorithm was bad, but because
the \emph{memory hierarchy} was being abused.

\section{The GPU memory hierarchy in one paragraph}

A modern data-centre GPU (H100, B200) has three relevant memory tiers:

\begin{itemize}
\item \term{HBM} (high-bandwidth memory): the on-package DRAM. 80 GB on
H100, ${\sim}3$ TB/s bandwidth.
\item \term{SRAM} (the on-chip shared memory): a few hundred KB per
streaming multiprocessor (SM), but \emph{much} faster — bandwidth on the
order of 19 TB/s.
\item Registers: even faster, smallest of all.
\end{itemize}

Going to HBM is roughly 6$\times$ slower than staying in SRAM. For an
operation that is bottlenecked by data movement rather than arithmetic,
choosing where to keep your tensors decides your wall clock.

\section{The standard implementation's mistake}

The naive way to implement attention on a GPU is what
equation~\eqref{eq:attn-mat} literally says:
\begin{enumerate}
\item Materialise the matrix \(S = QK\T \in \R^{n \times n}\) in HBM.
\item Apply a row-wise softmax to \(S\), writing \(A\) back to HBM.
\item Compute \(O = A V\) by reading \(A\) and \(V\) from HBM.
\end{enumerate}
For \(n = 64{,}000\), \(S\) is \({\sim}16\) GB. Writing it out and reading
it back twice is the bottleneck — not the arithmetic.

\section{What FlashAttention does}

The FlashAttention insight is a classical one from numerical linear algebra:
\term{tile} the computation, do it in pieces small enough to fit in SRAM,
and \emph{never write the intermediate \(n \times n\) matrix to HBM}.

The clever bit is doing softmax this way. Naively, softmax needs the full
row in order to compute the normalising constant
\(\sum_j \exp(s_{i,j})\). The standard trick is to compute it in two
passes — once to find the max, once to compute the exponentials and the
sum. \citet{dao2022flashattention} use a streaming, single-pass formulation
based on running statistics: as you process tile after tile of \(K, V\),
update running max and running sum in registers, and finalise the output for
each query tile when its work is done.

\section{What it costs and what it doesn't}

\paragraph{The good news.}
\begin{itemize}
\item \textbf{Memory.} Activation memory drops from \(\bigO(n^2)\) to
\(\bigO(n)\). At \(n = 64\)K, that is the difference between 16 GB and a
few MB.
\item \textbf{Wall clock.} 2--4$\times$ faster end-to-end for typical
sequence lengths in 2022; FlashAttention-2 \citep{dao2023fa2} and
FlashAttention-3 \citep{shah2024fa3} push this further by exploiting
GPU-specific instructions.
\item \textbf{Numerics.} The output is bit-equivalent to the naive
implementation (modulo standard floating-point reduction order). It is
\emph{exact} attention, not an approximation.
\end{itemize}

\paragraph{The bad news.}
\begin{equation}
  \text{FLOPs} \;=\; \bigO(n^2 d).
\end{equation}
Same as before. FlashAttention is a constant-factor win, not an asymptotic
one. Doubling \(n\) still quadruples the work; the work is just done much
more efficiently.

\keyidea{FlashAttention buys roughly an order of magnitude of practical
context length on the same hardware. It does not change the slope of the
cost curve. At sufficiently large \(n\), no implementation of dense
attention can be fast enough.}

\section{Why this matters for the rest of the book}

Two reasons.

\paragraph{It set the practical baseline.} ``Faster than FlashAttention-2 on
a B200'' is the de-facto unit of comparison for any new long-context
attention mechanism in 2026. When SubQ reports a 52$\times$ prefill
speedup at 1M tokens (Chapter~\ref{ch:benchmarks}), the baseline is
FlashAttention-2 — which is itself an aggressively-optimised implementation
of plain dense attention. Beating that by 52$\times$ is a real claim.

\paragraph{It clarified what \emph{can't} be fixed by implementation alone.}
Once your implementation is IO-optimal, what is left is the algorithm. If
the algorithm is \(n^2\), the algorithm has to change. That observation —
that we have collectively wrung most of the implementation lemon dry — is
why the architectural alternatives in Part~V get serious attention now,
when they were curiosities in 2020.

\begin{lstlisting}[caption={FlashAttention-1 idea: tiled streaming softmax so no full \(n\times n\) score matrix materialises in \textsc{hbm}; FA-2/FA-3 refine warp scheduling and Hopper overlaps (paragraph below).}]
# Outer loop: query tiles Q_i that fit in SRAM. Inner loop: K_j,V_j tiles.
O = zeros_like(V)                      # output [n, d_v], resident in HBM
for q_tile in range(0, n, T_q):
    Q_i = load(Q[q_tile:q_tile+T_q])   # HBM -> SRAM
    O_i = zeros(T_q, d_v); stats init   # registers/SRAM for softmax stats
    for kv_tile in range(0, n, T_k):
        K_j = load(K[kv_tile:kv_tile+T_k])   # HBM -> SRAM
        V_j = load(V[kv_tile:kv_tile+T_k])   # HBM -> SRAM
        S_ij = (Q_i @ K_j.T) * scale         # [T_q,T_k], stays in SRAM (never fully HBM)
        # Merge S_ij into running softmax + accumulate softmax(S_ij) @ V_j on-chip.
        # ...
        # (exact ``streaming softmax'' algebra keeps partial max/sum exp and rescales O_i)
    store(O[q_tile:q_tile+T_q], normalise(O_i))  # single write-back per query tile
\end{lstlisting}

\noindent
\citet{dao2022flashattention} spell out the exact streaming-softmax algebra used
inside the inner loop so partial \(S_{ij}\) tiles fuse with prior tiles without
ever forming the full \(n\times n\) map in \textsc{hbm}.

\paragraph{How FA-2 and FA-3 differ from FA-1.}
FlashAttention-2 \citep{dao2023fa2} keeps the same numerical streaming-softmax
recipe but changes \emph{work scheduling}: sequence-length parallelism warps are
assigned more evenly, recomputation for backward passes is tightened, and fewer
non-matmul instructions sit on the critical path, yielding roughly
\(1.5\)--\(2\times\) higher throughput on typical shapes without touching the
\(n^{2}\) flop count.

FlashAttention-3 \citep{shah2024fa3} targets Hopper (\texttt{fp8} tensor cores,
asynchronous \texttt{cp.async} copies, and warp-specialised softmax) so that
memory movement and instruction issue overlap more aggressively; again the
asymptotics stay quadratic, but wall-clock drops another \(\mathcal{O}(1)\)
factor on hardware where FA-1 was already close to roofline.

Honest bottom line: each generation bought roughly a \(1.5\)--\(2\times\) win on
real benchmarks when kernels were bound by \textsc{hbm} bandwidth or occupancy,
not a change in \(\Theta(n^{2})\) scaling---that change remains the province of
Part~V architectures.
